\documentclass{article}
\usepackage{amsmath, amssymb, amsthm}
\usepackage{hyperref}
\usepackage{geometry}
\geometry{margin=1in}

\title{Erd\H{o}s Problem \#271}
\date{Last updated: 2025-08-31}
\author{Source: erdosproblems.com}

\begin{document}
\maketitle

\noindent\textbf{Status:} OPEN \quad \textbf{No prize}

\noindent\textbf{Tags:} additive combinatorics, arithmetic progressions

\medskip
\noindent\textbf{Note:} Stanley sequences

\medskip

\noindent\textbf{Problem Statement:}

Let $A(n)=\{a_0&#60;a_1&#60;\cdots\}$ be the sequence defined by $a_0=0$ and $a_1=n$, and for $k\geq 1$ define $a_{k+1}$ as the least positive integer such that there is no three-term arithmetic progression in $\{a_0,\ldots,a_{k+1}\}$. Can the $a_k$ be explicitly determined? How fast do they grow?




It is easy to see that $A(1)$ is the set of integers which have no 2 in their base 3 expansion. Odlyzko and Stanley \cite{OdSt78} have found similar characterisations are known for $A(3^k)$ and $A(2\cdot 3^k)$ for any $k\geq 0$ and conjectured in general that such a sequence always eventually either satisfies\[a_k\asymp k^{\log_23}\]or\[a_k \asymp \frac{k^2}{\log k}.\]There is no known sequence which satisfies the second growth rate, but Lindhurst \cite{Li90} gives data which suggests that $A(4)$ has such growth ($A(4)$ is given as A005487 in the OEIS). Moy \cite{Mo11} has proved that, for all such sequences, for all $\epsilon&gt;0$, $a_k\leq (\frac{1}{2}+\epsilon)k^2$ for all sufficiently large $k$. van Doorn and Sothanaphan have noted in the comment section that Moy's proof can be upgraded to give a fully explicit result of\[a_k\leq \frac{(k-1)(k+2)}{2}+n\]for all $k\geq 0$. In general, sequences which begin with some initial segment and thereafter are continued in a greedy fashion to avoid three-term arithmetic progressions are known as Stanley sequences.




References

[Li90] S. Lindhurst, An investigation of several interesting sets of numbers generated by the greedy algorithm . Senior thesis at Princeton University (1990).

[Mo11] Moy, Richard A., On the growth of the counting function of Stanley sequences . Discrete Math. (2011), 560-562.

[OdSt78] A. Odlyzko and R. Stanley, Some curious sequences constructed with the greedy algorithm . Bell Laboratories internal memorandum (1978).

\noindent\textbf{Related OEIS sequences:} A005487


\bigskip
\noindent\small{Source: \url{https://www.erdosproblems.com/271}}
\end{document}
